%% maestro2tex -- generated table; do not edit by hand.
\begin{table}[htbp]
  \centering
  \caption[{DC accuracy by process corner.}]{DC accuracy by process corner. Closed transimpedance loop bench (Figure~\ref{fig:tb-tiafb}): $R_f = \SI{560}{\kilo\ohm}$, $C_f = \SI{22}{\pico\farad}$, $C_L = \SI{5}{\pico\farad}$, $C_c = \SI{500}{\femto\farad}$, \texttt{BiasAdj}~=~37, $v_{\mathrm{cm}} = \SI{1.25}{\volt}$, \SI{2.5}{\volt} supply, Randles electrode with $R_s = \SI{100}{\ohm}$, $R_{\mathrm{ct}} = \SI{1}{\mega\ohm}$, $C_{\mathrm{dl}} = \SI{100}{\nano\farad}$. Point \texttt{tt} is the typical statistical point at \SI{37}{\celsius}; the four skew corners are at \SI{25}{\celsius}. Gain error and INL are evaluated over the central \(\pm 0.9\) of the full-scale input-current range and are referred to the nominal \SI{2.441}{\milli\volt} ADC LSB; the converter's measured step is \SI{1.591}{\milli\volt} (Table~\ref{tab:adc-transfer}), so multiply the LSB row by \num{1.535} to read it as converter codes. These are corner figures with zero statistical offset: the offset a manufactured channel carries is the Monte Carlo result of Table~\ref{tab:tiaab-mc-offset}. These corner spreads are a lower bound: the corner set holds the passive model rows at typical (pinned-passive caveat, page~\pageref{p:pinned-passives}).}
  \label{tab:tiaab-accuracy}
  \begin{tabular}{lrrrrrrr}
    \toprule
    Parameter & Unit & tt & ff & fs & sf & ss & Spread \\
    \midrule
    Output offset at zero current & \si{\milli\volt} & $-0.967$ & $-1.056$ & $-0.987$ & $-0.832$ & $-0.762$ & 0.294 \\
    Gain error & ppm & $-167.3$ & $-221.9$ & $-150.6$ & $-161.0$ & $-106.3$ & 115.6 \\
    Integral non-linearity & LSB & 0.021 & 0.033 & 0.016 & 0.023 & 0.011 & 0.022 \\
    \bottomrule
  \end{tabular}
\end{table}
